\section{Cost and Token Efficiency in Agent-Based Workflows}

An agent's cost depends more on how much context builds up over a task than on any single request. Each step in an agent's work resends the conversation so far, the tool definitions it has access to, and any project context it was given, so the longer or more complex a task gets, the more expensive every additional step becomes --- and generated output typically costs several times more per token than the input, so unnecessarily long or repetitive agent output adds up faster than it might seem.

\textbf{A first practice is reusing stable, repeated context instead of resending it in full every time.} Part 2's Charter and structure maps are a concrete instance of this: instead of retyping the project's architectural rules and conventions into every task prompt, they were kept as standing files that an agent read at the start of a task. This avoided resending the same rules and conventions in full each time a new feature was requested, and is usually the easiest cost saving to put in place relative to the effort it takes.

\textbf{A second practice is matching each task to the model it actually needs, instead of using the most capable model for everything.} Simple, well-defined subtasks can be handled by a smaller or cheaper model, while a stronger model is reserved for steps that genuinely require deeper reasoning. This mirrors the same idea behind splitting work across specialized agent roles, just applied to choosing a model instead of dividing up a task.

\textbf{A third practice is being deliberate about what context is actually included in each step, rather than letting it grow by default.} Part 2's Charter applies this directly to the Reviewer Agent: rather than exploring the broader codebase, it is restricted to reading exactly three named artifacts --- the approved implementation plan, the walkthrough, and the specific files listed as changed --- with the Charter stating this restriction exists explicitly to keep the review "focused and token-efficient." Pulling in only the specific information a step needs, instead of full documents or the entire codebase every time, keeps cost from creeping up simply because it was easier not to trim it.

